\documentclass[11pt]{article}

\usepackage[T1]{fontenc}
\usepackage[utf8]{inputenc}
\usepackage{lmodern}
\usepackage[margin=1in]{geometry}
\usepackage{amsmath,amssymb,amsthm,mathtools}
\usepackage{booktabs,array}
\usepackage[shortlabels]{enumitem}
\usepackage{xcolor}
\usepackage{hyperref}
\usepackage{xurl}
\usepackage{microtype}
\usepackage[numbers,sort&compress]{natbib}

\hypersetup{colorlinks=true,linkcolor=blue,citecolor=blue,urlcolor=blue}
\setlength{\parindent}{0pt}
\setlength{\parskip}{0.55em}
\setlength{\emergencystretch}{3em}
\renewcommand{\arraystretch}{1.12}
\newcolumntype{L}[1]{>{\raggedright\arraybackslash}p{#1}}

% Permit line breaks after underscores inside typewriter identifiers.
\renewcommand{\_}{\textunderscore\allowbreak}

\newtheorem{theorem}{Theorem}
\newtheorem{definition}{Definition}
\newtheorem{proposition}{Proposition}
\theoremstyle{remark}
\newtheorem{remark}{Remark}

\newcommand{\C}{\mathbb C}
\newcommand{\R}{\mathbb R}
\newcommand{\Mn}{\mathrm{M}_n(\C)}
\newcommand{\Tr}{\operatorname{Tr}}
\newcommand{\norm}[1]{\lVert #1\rVert}
\newcommand{\one}{\mathbf 1}
\newcommand{\pinch}{\mathcal E}
\newcommand{\lued}{\mathcal L}
\newcommand{\ip}[2]{\langle #1,\,#2\rangle}

\title{\textbf{Machine-Checked Finite Quantum Event Algebras:\\
Born Weights, L\"uders Conditioning, Central Conditional Expectations,\\
and the Tsirelson Bound in Lean~4}}
\author{Bernhard Mueller\thanks{Corresponding author:
\href{mailto:bernhard@floatingpragma.ai}{\texttt{bernhard@floatingpragma.ai}}}\\
Pragma Research Inc., Washington, United States\\
\href{mailto:bernhard@floatingpragma.ai}{\texttt{bernhard@floatingpragma.ai}}}
\date{}

\begin{document}
\maketitle

\begin{abstract}
We present a machine-checked development of the elementary theory of finite-dimensional quantum event algebras: events as Hermitian idempotents in the algebra of $n \times n$ complex matrices, states as positive-semidefinite matrices of unit trace, Born weights as the trace pairing, L\"uders conditioning with its repeatability and fixed-point structure, the conditional expectation onto the commutative subalgebra generated by a projective partition of unity, the positive expectation functional of a state, and the Tsirelson bound for the CHSH combination. The development comprises 64 lemmas and 8 definitions in Lean~4 over the Mathlib library, with no proof placeholders. An axiom audit embedded in the source and re-executed on every build reports, for each result, only the three standard axioms of Lean's classical mathematics; the core CHSH ring identity avoids the axiom of choice entirely. Every lemma is tagged as provable from the $*$-algebra axioms alone or as consuming a declared tracial state, so the boundary between algebraic and probabilistic content is machine-visible. The Tsirelson bound is proved in operator-norm form in an abstract unital C*-setting that assumes neither completeness nor an order structure, then instantiated to finite matrix algebras; it complements the order-form inequality already present in Mathlib and consumes Mathlib's CHSH hypothesis interface for interoperability. We describe the formalization module by module, document the Mathlib gaps and workarounds encountered, and situate the development relative to the Isabelle/HOL line on projective measurements and the CHSH inequality.
\end{abstract}

\noindent\textbf{Keywords:} quantum probability, event algebras, Born rule, L\"uders conditioning, conditional expectation, Tsirelson bound, formalized mathematics, Lean~4, Mathlib

\medskip
\noindent\textbf{Mathematics Subject Classification (2020):} Primary 81P15, 68V20; Secondary 81P40, 46L53, 81P10.

\section{Introduction}
\label{sec:intro}

Finite-dimensional event algebras support a complete and elementary account of quantum probability. Over the matrix algebra $\Mn$, events are orthogonal projections, states are density matrices, the Born weight of an event under a state is a trace pairing, conditioning is the L\"uders update \cite{lueders1950,buschlahti2009}, classicality is captured by commutative subalgebras and conditional expectations onto them, and the characteristic nonclassical phenomenon, the violation of the CHSH inequality up to Tsirelson's bound \cite{bell1964,chsh1969,tsirelson1980}, is already visible for $4 \times 4$ matrices. Every object in this list is finite-dimensional linear algebra; no unbounded operators, no measure theory, and no functional-analytic subtleties intervene. This elementary core is also where much of the current foundational discussion lives: the status and derivation of the Born rule \cite{reddigerpoirier2025,aharonovshushi2024}, the relation between quantum phenomena and Kolmogorov probability \cite{reddiger2026}, and the derivability of the von Neumann--L\"uders measurement postulate \cite{goldstein2026} remain active topics, including in this journal.

Published treatments of this core, from von Neumann \cite{vonneumann1932} onward, leave the bookkeeping informal. Which statements require positivity of the state, which require only hermiticity, and which are identities of the trace calculus valid for arbitrary matrices? What exactly happens when one conditions on an event of weight zero? Is orthogonality of projections a symmetric hypothesis or a symmetric conclusion? Where, precisely, does the trace enter, and where does pure $*$-algebra suffice? In informal mathematics these questions are resolved silently, case by case, by the reader. A proof assistant forces every such choice into the open and then certifies the result.

This paper presents the first Lean~4/Mathlib formalization of the finite quantum event-algebra core organized around a distinguished commutative center and an explicit conditional-expectation projector, with Born weights, repeatable L\"uders conditioning, and a finite CHSH/Tsirelson bound, under an audited axiom base and an explicit separation between results proved from the algebraic axioms alone and results consuming a declared tracial state. Concretely, the development delivers:
\begin{enumerate}[(a)]
\item 64 lemmas and 8 definitions, in Lean~4 \cite{lean4} over Mathlib \cite{mathlib2020}, with no \texttt{sorry} placeholders and no axioms beyond Lean's standard classical base;
\item an axiom audit embedded in the source: every lemma is followed by a \texttt{\#print axioms} command, and each such audit reports exactly \texttt{[propext, Classical.choice, Quot.sound]}; the CHSH ring identity \texttt{chsh\_mul\_self} reports only \texttt{[propext, Quot.sound]}, avoiding the axiom of choice;
\item a tagging discipline under which every lemma is classified in its documentation string as \emph{algebra-only} (content provable in the $*$-algebra, with norm structure where relevant) or \emph{consumes a tracial state} (content passing through the trace pairing $(\rho,P) \mapsto \Tr(\rho P)$);
\item the L\"uders update packaged with repeatability, idempotence, composition on commuting events, the classical restriction, and a fixed-point characterization identifying the fixed states of conditioning with the states certain of the event;
\item the pinching map onto the commutant of a projective partition of unity, verified to satisfy the full conditional-expectation package: projector, state-preserving, selfadjoint for the trace inner product, Pythagorean, contractive, trace-compatible, unique, and commuting with L\"uders conditioning on central events;
\item the Tsirelson bound in operator-norm form, proved in an abstract unital C*-setting that assumes neither completeness nor any order structure, with the exact ring identity underlying it isolated as a standalone lemma, an instantiation to finite matrix algebras, and an interface lemma consuming Mathlib's existing CHSH hypothesis bundle.
\end{enumerate}
The development builds with the toolchain \texttt{leanprover/lean4:v4.29.1} and the Mathlib version pinned to match; it imports Mathlib and nothing else.

Three scope limitations should be stated at the outset. First, the novelty claim above is deliberately narrow and organizational. Formalizations of projective measurement, the Born rule for projective effects, and the CHSH/Tsirelson inequality exist in the Isabelle/HOL ecosystem, in particular in the work of Echenim and Mhalla \cite{echenim2021afp,echenimmhalla2023afp,echenimmhalla2024}; we make no claim of first formalization for any of those individual results. The contribution here is the organization of the finite core around a distinguished commutative center with an explicit conditional-expectation projector, the lemma-by-lemma separation of algebraic from state-consuming content, the abstract C*-setting of the norm-form bound, and the prover ecosystem. Second, no attainment or tightness statement for the Tsirelson bound is formalized: the development ships the inequality only. Third, everything is finite-dimensional; extensions toward infinite dimensions are discussed as outlook in Section~\ref{sec:conclusion}.

Finite event algebras arise, for example, as the measurement cores of finite models of quantum theory. Beyond that single remark, this paper is self-contained: Section~\ref{sec:setting} fixes the mathematical setting, Section~\ref{sec:arch} describes the architecture and conventions of the formalization, Sections~\ref{sec:basic} through \ref{sec:tsirelson} walk through the five modules with the main statements as numbered propositions carrying their Lean names, Section~\ref{sec:gaps} documents Mathlib gaps and proof engineering, Section~\ref{sec:related} discusses related work, and Section~\ref{sec:conclusion} concludes.

\section{Mathematical setting}
\label{sec:setting}

Throughout, $n$ and $k$ denote natural numbers and $\Mn$ denotes the $*$-algebra of $n \times n$ complex matrices, with conjugate transpose $X \mapsto X^{*}$, trace $\Tr$, identity $\one$, and zero $0$. All statements below are formalized for every $n$, including the degenerate case $n = 0$; only the matrix instantiation of the Tsirelson bound (Proposition~\ref{prop:tsirelson-instances}) assumes $n \neq 0$.

The partial order on $\C$ is the one in which $z \le w$ exactly when $\operatorname{Re} z \le \operatorname{Re} w$ and $\operatorname{Im} z = \operatorname{Im} w$; in particular $0 \le z$ asserts that $z$ is real and nonnegative. Nonnegativity and boundedness of Born weights are proved below in this partial order, which is strictly stronger than the corresponding assertions about real parts; real-part corollaries are provided alongside each such statement.

\begin{definition}[Event]
\label{def:event}
An \emph{event} of $\Mn$ is a Hermitian idempotent: a matrix $P \in \Mn$ with $P^{*} = P$ and $P^2 = P$. Lean name: \texttt{IsEvent}.
\end{definition}

Events are exactly the orthogonal projections, the sharp yes/no tests of the finite quantum formalism \cite{vonneumann1932}. The \emph{subevent} relation is witnessed algebraically: $P$ is a subevent of $Q$ when $PQ = P$.

\begin{definition}[State]
\label{def:state}
A \emph{state} on $\Mn$ is a density matrix: a positive-semidefinite $\rho \in \Mn$ with $\Tr \rho = 1$. Lean name: \texttt{IsState}.
\end{definition}

\begin{definition}[Born weight]
\label{def:bornweight}
The \emph{Born weight} of $P$ under $\rho$ is the trace pairing
\[
\mu_\rho(P) \coloneqq \Tr(\rho P) \in \C .
\]
Lean name: \texttt{bornWeight}.
\end{definition}

The weight is defined as a complex number for arbitrary matrix arguments; reality under hermiticity is a theorem (Proposition~\ref{prop:reality}), and the probabilistic bounds hold under the appropriate positivity hypotheses (Proposition~\ref{prop:prob}).

\begin{definition}[L\"uders update]
\label{def:lueders}
The \emph{L\"uders update} of $\rho$ by $P$ is the compressed and renormalized matrix
\[
\lued_P(\rho) \coloneqq \bigl(\mu_\rho(P)\bigr)^{-1}\, P \rho P .
\]
Lean name: \texttt{luedersUpdate}.
\end{definition}

The scalar inverse is taken in $\C$ with Lean's convention $0^{-1} = 0$, so conditioning on an event of vanishing weight degenerates to the zero matrix; every substantive lemma about the update carries the explicit guard $\mu_\rho(P) \neq 0$. The rule goes back to L\"uders \cite{lueders1950}; see Busch and Lahti \cite{buschlahti2009} for its place in the measurement literature.

\begin{definition}[Center partition]
\label{def:partition}
A \emph{center partition} of $\Mn$ is a finite family $\pi = (P_i)_{i < k}$ of events that are pairwise orthogonal, $P_i P_j = 0$ for $i \neq j$, and complete, $\sum_{i < k} P_i = \one$. Lean name: \texttt{CenterPartition}.
\end{definition}

The members of a center partition commute pairwise (Proposition~\ref{prop:partition-alg}), so the partition generates a commutative unital subalgebra of $\Mn$: the distinguished \emph{center} of the measurement context. Its \emph{commutant} is
\[
\pi' \coloneqq \{\, X \in \Mn : X P_i = P_i X \text{ for all } i < k \,\}.
\]

\begin{definition}[Center expectation]
\label{def:pinching}
The \emph{center expectation} (pinching map) attached to a center partition $\pi$ is
\[
\pinch_\pi(X) \coloneqq \sum_{i < k} P_i X P_i .
\]
Lean name: \texttt{centerExpectation}.
\end{definition}

Section~\ref{sec:center} verifies that $\pinch_\pi$ satisfies, in the finite quantum setting, the same package of properties that characterizes a classical conditional expectation: it is a projector onto $\pi'$, maps states to states, is selfadjoint for the trace inner product $\ip{X}{Y} = \Tr(X^{*}Y)$, satisfies the Pythagorean identity, contracts the squared trace norm, is uniquely determined by trace compatibility, and commutes with L\"uders conditioning on central events.

Finally, the CHSH combination. In a unital C*-algebra $A$, take selfadjoint involutions $a_0, a_1, b_0, b_1$ (each equal to its own star, each squaring to $\one$) such that every $a_i$ commutes with every $b_j$, and form
\[
S \coloneqq a_0 b_0 + a_0 b_1 + a_1 b_0 - a_1 b_1 .
\]
Classically, dichotomic $\pm 1$ observables obey $\lvert \langle S \rangle \rvert \le 2$ \cite{bell1964,chsh1969}; quantum mechanically the sharp bound is Tsirelson's $\norm{S} \le 2\sqrt{2}$ \cite{tsirelson1980}. The bridge from the event layer to these hypotheses is the dichotomic observable $\one - 2P$ of an event $P$, which is a selfadjoint involution (Proposition~\ref{prop:closure}(ix)).

\section{The formalization: architecture and conventions}
\label{sec:arch}

\subsection{Modules}

The development consists of five modules under a single namespace, importing Mathlib and nothing else:

\begin{table}[ht]
\centering
\small
\begin{tabular}{@{}lL{4.9cm}L{4.6cm}c@{}}
\toprule
Module & Content & Definitions & Lemmas\\
\midrule
\texttt{Basic} & events, states, Born weights & \texttt{IsEvent}, \texttt{IsState}, \texttt{bornWeight} & 25\\
\texttt{Lueders} & conditioning, certainty sets, fixed points & \texttt{luedersUpdate}, \texttt{certainStates} & 12\\
\texttt{CenterExpectation} & partitions, pinching, conditional expectation & \texttt{CenterPartition}, \texttt{centerExpectation} & 16\\
\texttt{StateFromTrace} & the expectation functional of a state & \texttt{expectation} & 5\\
\texttt{Tsirelson} & CHSH ring identity, norm-form Tsirelson bound & (none) & 6\\
\midrule
Total & & 8 & 64\\
\bottomrule
\end{tabular}
\caption{Module layout of the development.}
\label{tab:modules}
\end{table}

The event and state layers are formalized over the concrete carrier \texttt{Matrix (Fin n) (Fin n) $\C$}. This choice trades generality for the density of Mathlib's finite matrix API (positive semidefiniteness, trace calculus, spectral theorem), which carries most of the proofs. The Tsirelson module, whose content is norm-theoretic rather than matrix-theoretic, is instead stated over an abstract typeclass context and instantiated to matrices only at the end (Section~\ref{sec:tsirelson}).

\subsection{Scoped instances}

Three pieces of structure that the development relies on are \emph{scoped} in Mathlib rather than globally active, and must be opened explicitly:
\begin{itemize}
\item \texttt{ComplexOrder}: the partial order on $\C$ of Section~\ref{sec:setting}, together with the ordered-ring structure that makes products of nonnegative complex numbers nonnegative;
\item \texttt{MatrixOrder}: the Loewner order on matrices;
\item \texttt{Matrix.Norms.L2Operator}: the $L^2$ operator norm on $\Mn$ together with its normed-ring and C*-identity instances. There is no globally installed norm on the matrix type, so the matrix instantiation of the Tsirelson bound lives behind this scope.
\end{itemize}
Section~\ref{sec:gaps} discusses the practical consequences of this scoping.

\subsection{The algebra-only/tracial split}

Every lemma in the development carries one of two tags in its documentation string:
\begin{itemize}
\item \textbf{algebra-only}: the content is pure $*$-algebra (ring and star structure; in the Tsirelson module, also norm structure), consuming no state and no trace functional;
\item \textbf{consumes a tracial state}: the content passes through the trace pairing $(\rho, P) \mapsto \Tr(\rho P)$.
\end{itemize}
The split is a statement-classification discipline rather than an import boundary: it makes visible, lemma by lemma, exactly where probabilistic content enters the theory. All event-closure lemmas, the entire partition algebra, the projector laws of the pinching map, and the entire Tsirelson module are algebra-only; every statement about Born weights, updates of states, and trace inner products consumes the trace. In the propositions below, this classification is reproduced: statements in Sections~\ref{sec:basic} through \ref{sec:functional} are tagged individually, and everything in Section~\ref{sec:tsirelson} is algebra-only.

\subsection{Axiom audit and build discipline}

Each of the 64 lemmas is followed in the source by a \texttt{\#print axioms} command, so the axiom audit is part of the source text and re-executes on every build. Every audit reports exactly
\[
\texttt{[propext, Classical.choice, Quot.sound]},
\]
the three standard axioms of Lean's classical mathematics (propositional extensionality, choice, and quotient soundness). The ring identity \texttt{chsh\_mul\_self} reports only \texttt{[propext, Quot.sound]}: the algebraic heart of the Tsirelson bound is choice-free. The development contains no \texttt{sorry} placeholders, no custom axioms, and no unfinished obligations; the toolchain is \texttt{leanprover/lean4:v4.29.1} with Mathlib pinned at the matching version, and the project builds cleanly under \texttt{lake} (8254 jobs).

Proofs in this paper are given as sketches only; the complete formal proofs constitute the accompanying development, and each proposition cites the corresponding Lean names in \texttt{typewriter} font.

\section{Events and Born weights}
\label{sec:basic}

The first module establishes closure of the event set under the natural logical operations and the probabilistic behavior of the Born weight.

\begin{proposition}[Event closure; algebra-only]
\label{prop:closure}
Let $P, Q \in \Mn$ be events. Then:
\begin{enumerate}[(i)]
\item $0$ and $\one$ are events (\texttt{isEvent\_zero}, \texttt{isEvent\_one});
\item the complement $\one - P$ is an event (\texttt{IsEvent.compl});
\item orthogonality is symmetric: if $PQ = 0$ then $QP = 0$ (\texttt{IsEvent.orthogonal\_symm});
\item if $PQ = 0$ then $P + Q$ is an event (\texttt{IsEvent.add});
\item if $PQ = QP$ then $PQ$ is an event (\texttt{IsEvent.mul\_of\_commute});
\item subevent absorption: if $PQ = P$ then $QP = P$ (\texttt{IsEvent.absorb\_of\_le});
\item if $PQ = P$ then $Q - P$ is an event (\texttt{IsEvent.sub\_of\_le});
\item every event is positive semidefinite (\texttt{IsEvent.posSemidef});
\item $\one - 2P$ is a selfadjoint involution: $(\one - 2P)^{*} = \one - 2P$ and $(\one - 2P)^2 = \one$ (\texttt{IsEvent.one\_sub\_two\_smul\_involution}).
\end{enumerate}
\end{proposition}

\begin{proof}[Proof sketch]
All items are short computations in the $*$-algebra. Symmetry of orthogonality (iii) and absorption (vi) take the conjugate transpose of the hypothesis and use hermiticity of both matrices; positive semidefiniteness (viii) exhibits $P = P^{*}P$ as a Gram matrix.
\end{proof}

Item (ix) is the bridge from events to the CHSH hypotheses of Section~\ref{sec:tsirelson}: each event answers a yes/no question, and $\one - 2P$ is the associated $\pm 1$-valued observable.

\begin{proposition}[Trace-pairing calculus; consumes a tracial state]
\label{prop:pairing}
For all $\rho, P, Q \in \Mn$, all finite families $(P_i)_{i \in s}$, and all $c \in \C$:
\begin{enumerate}[(i)]
\item $\mu_\rho(P + Q) = \mu_\rho(P) + \mu_\rho(Q)$ and $\mu_\rho(P - Q) = \mu_\rho(P) - \mu_\rho(Q)$ (\texttt{bornWeight\_add}, \texttt{bornWeight\_sub});
\item $\mu_\rho\bigl(\sum_{i \in s} P_i\bigr) = \sum_{i \in s} \mu_\rho(P_i)$ (\texttt{bornWeight\_sum});
\item $\mu_{c\rho}(P) = c\, \mu_\rho(P)$ (\texttt{bornWeight\_smul});
\item if $P^2 = P$ then $\Tr(P \rho P) = \mu_\rho(P)$ (\texttt{trace\_sandwich}).
\end{enumerate}
\end{proposition}

Item (iv), the sandwich identity, requires only idempotence of $P$; hermiticity is never used. It is the workhorse that lets positivity arguments compress the state by the event and is applied throughout the conditioning module.

\begin{proposition}[Reality; consumes a tracial state]
\label{prop:reality}
If $\rho$ and $P$ are Hermitian (neither positivity nor normalization is assumed), then $\overline{\mu_\rho(P)} = \mu_\rho(P)$; equivalently, the weight equals its own real part (\texttt{star\_bornWeight}, \texttt{bornWeight\_eq\_re}).
\end{proposition}

\begin{proposition}[Probabilistic bounds; consumes a tracial state]
\label{prop:prob}
Let $P, Q$ be events. Then, in the partial order of $\C$:
\begin{enumerate}[(i)]
\item if $\rho$ is positive semidefinite, then $0 \le \mu_\rho(P)$ (\texttt{bornWeight\_nonneg}), with real-part form \texttt{bornWeight\_re\_nonneg};
\item if $\rho$ is a state, then $\mu_\rho(\one) = 1$ (\texttt{bornWeight\_one});
\item if $PQ = 0$, then $\mu_\rho(P + Q) = \mu_\rho(P) + \mu_\rho(Q)$ (\texttt{bornWeight\_add\_of\_orthogonal});
\item if $\rho$ is a state, then $\mu_\rho(P) \le 1$ (\texttt{bornWeight\_le\_one}), with real-part form \texttt{bornWeight\_re\_le\_one};
\item if $\rho$ is positive semidefinite and $PQ = P$, then $\mu_\rho(P) \le \mu_\rho(Q)$ (\texttt{bornWeight\_mono});
\item if $\rho$ is positive semidefinite, then $\mu_\rho(P) \neq 0$ exactly when $0 < \operatorname{Re} \mu_\rho(P)$ (\texttt{bornWeight\_ne\_zero\_iff\_re\_pos}).
\end{enumerate}
\end{proposition}

\begin{proof}[Proof sketch]
For (i), compress: $\mu_\rho(P) = \Tr(P \rho P)$ by Proposition~\ref{prop:pairing}(iv), and $P \rho P = P \rho P^{*}$ is a positive-semidefinite sandwich, whose trace is nonnegative. The upper bound (iv) applies (i) to the complement $\one - P$; monotonicity (v) applies (i) to the difference event $Q - P$ of Proposition~\ref{prop:closure}(vii).
\end{proof}

Together, Propositions~\ref{prop:closure} and \ref{prop:prob} verify that $P \mapsto \mu_\rho(P)$ is a finitely additive, normalized, monotone probability assignment on events, with reality and the bounds stated in the strong ($\C$-order) form. Item (vi) is the exact interface later needed by the conditioning guards.

\section{L\"uders conditioning}
\label{sec:lueders}

The second module packages the L\"uders update as a state transformer and proves its characteristic laws. A remark on hypothesis minimality is in order: because the formal statements register exactly what each proof uses, several conditioning identities below hold for an \emph{arbitrary} matrix $\rho$, with state hypotheses appearing only where positivity or normalization genuinely enters.

\begin{proposition}[State property; consumes a tracial state]
\label{prop:update-state}
Let $P$ be an event.
\begin{enumerate}[(i)]
\item If $\rho$ is positive semidefinite, then $\lued_P(\rho)$ is positive semidefinite, with no guard: in the degenerate case the update is the zero matrix (\texttt{luedersUpdate\_posSemidef}).
\item If $\mu_\rho(P) \neq 0$, then $\Tr \lued_P(\rho) = 1$ for arbitrary $\rho$ (\texttt{trace\_luedersUpdate}).
\item If $\rho$ is a state and $\mu_\rho(P) \neq 0$, then $\lued_P(\rho)$ is a state (\texttt{luedersUpdate\_isState}).
\end{enumerate}
\end{proposition}

By Proposition~\ref{prop:prob}(vi), the guard $\mu_\rho(P) \neq 0$ is equivalent, for positive-semidefinite $\rho$, to the weight having strictly positive real part.

\begin{proposition}[Repeatability; consumes a tracial state]
\label{prop:repeat}
Let $P$ be an event and $\mu_\rho(P) \neq 0$. Then
\[
\mu_{\lued_P(\rho)}(P) = 1
\qquad (\texttt{bornWeight\_luedersUpdate\_self}).
\]
\end{proposition}

Neither positivity nor normalization of $\rho$ is used: repeatability is an identity of the trace calculus, valid for every matrix $\rho$ subject to the guard.

\begin{proposition}[Idempotence; consumes a tracial state]
\label{prop:idem}
Let $P$ be an event and $\mu_\rho(P) \neq 0$. Then
\[
\lued_P\bigl(\lued_P(\rho)\bigr) = \lued_P(\rho)
\qquad (\texttt{luedersUpdate\_idem}).
\]
\end{proposition}

\begin{proposition}[Commuting composition and order exchange; consumes a tracial state]
\label{prop:compose}
Let $P, Q$ be events with $PQ = QP$.
\begin{enumerate}[(i)]
\item If $\mu_\rho(P) \neq 0$, then $\lued_Q\bigl(\lued_P(\rho)\bigr) = \lued_{PQ}(\rho)$ (\texttt{luedersUpdate\_luedersUpdate\_of\_commute}); the product $PQ$ is an event by Proposition~\ref{prop:closure}(v). Only the first guard is needed: if the joint weight vanishes, both sides degenerate to zero together.
\item If moreover $\mu_\rho(Q) \neq 0$, then conditioning is order-exchangeable: $\lued_Q\bigl(\lued_P(\rho)\bigr) = \lued_P\bigl(\lued_Q(\rho)\bigr)$ (\texttt{luedersUpdate\_comm}).
\end{enumerate}
\end{proposition}

\begin{proposition}[Classical restriction; consumes a tracial state]
\label{prop:classical}
Let $P$ be an event and suppose the state matrix commutes with it, $\rho P = P \rho$. Then
\[
\lued_P(\rho) = \bigl(\mu_\rho(P)\bigr)^{-1}\, \rho P
\qquad (\texttt{luedersUpdate\_of\_commute}),
\]
with no nonvanishing guard: the identity holds for arbitrary $\rho$ under the stated commutation.
\end{proposition}

In the commuting case the L\"uders update is the normalized restriction, which is the formula of classical conditional probability; quantum conditioning strictly extends the classical rule.

\subsection{Conditioning as one-step convergence to a fixed point}

Repeatability and idempotence assemble into a fixed-point description of measurement update.

\begin{definition}[Certainty set]
\label{def:certain}
The \emph{certainty set} of an event $P$ is
\[
\mathcal C_P \coloneqq \{\, \sigma \in \Mn : \sigma \text{ is a state and } \mu_\sigma(P) = 1 \,\}.
\]
Lean name: \texttt{certainStates}.
\end{definition}

\begin{proposition}[Fixed-point characterization; consumes a tracial state]
\label{prop:fixed}
Let $P$ be an event.
\begin{enumerate}[(i)]
\item If $\rho$ is a state with $\mu_\rho(P) \neq 0$, then $\lued_P(\rho) \in \mathcal C_P$: conditioning lands in the certainty set in a single step (\texttt{luedersUpdate\_mem\_certainStates}).
\item Support: if $\sigma \in \mathcal C_P$, then $\sigma P = \sigma = P \sigma$ (\texttt{mul\_eq\_self\_of\_bornWeight\_one}).
\item Every $\sigma \in \mathcal C_P$ is a fixed point: $\lued_P(\sigma) = \sigma$ (\texttt{luedersUpdate\_eq\_self\_of\_mem\_certainStates}).
\item Conversely, for a state $\sigma$ with $\mu_\sigma(P) \neq 0$,
\[
\lued_P(\sigma) = \sigma \iff \sigma \in \mathcal C_P
\qquad (\texttt{luedersUpdate\_eq\_self\_iff}).
\]
\end{enumerate}
\end{proposition}

\begin{proof}[Proof sketch]
The substantive item is (ii). The complement $Q = \one - P$ has weight zero under $\sigma$, so the compression $Q \sigma Q$ is positive semidefinite with zero trace and hence vanishes. The Cauchy--Schwarz inequality for the quadratic form of $\sigma$ (in the form: if $\langle x, \sigma x\rangle = 0$ for a positive-semidefinite $\sigma$ then $\sigma x = 0$) then kills $\sigma Q$ itself, column by column. Items (iii) and (iv) unfold the update using (ii).
\end{proof}

\begin{remark}
Among the states assigning $P$ nonzero weight, conditioning on $P$ is a retraction onto its own fixed-point set $\mathcal C_P$: measurement update is convergence to a fixed point, reached after one step. The formalization thus supports the reading of the L\"uders rule as a projection onto certainty rather than a dynamical process with transient behavior.
\end{remark}

\section{Conditional expectation onto a commutative center}
\label{sec:center}

The third module is the organizational core of the development. A center partition $\pi = (P_i)_{i<k}$ (Definition~\ref{def:partition}) plays the role of a classical measurement context; the pinching map $\pinch_\pi$ (Definition~\ref{def:pinching}) is verified to be \emph{the} conditional expectation onto its commutant, in the precise sense of the uniqueness statement below.

\begin{proposition}[Partition algebra; algebra-only]
\label{prop:partition-alg}
For a center partition $\pi$ and $i, j < k$:
\begin{enumerate}[(i)]
\item $P_i P_j = \delta_{ij} P_i$ (\texttt{CenterPartition.proj\_mul\_proj});
\item $P_i P_j = P_j P_i$: the partition generates a commutative subalgebra (\texttt{CenterPartition.proj\_commute}).
\end{enumerate}
\end{proposition}

\begin{proposition}[Projector laws; algebra-only]
\label{prop:projector}
For every $X \in \Mn$ and $i < k$:
\begin{enumerate}[(i)]
\item absorption: $\pinch_\pi(X)\, P_i = P_i X P_i = P_i\, \pinch_\pi(X)$ (\texttt{centerExpectation\_mul\_proj}, \texttt{proj\_mul\_centerExpectation});
\item $\pinch_\pi(X) \in \pi'$: the pinching is commutant-valued (\texttt{proj\_commute\_centerExpectation});
\item $\pinch_\pi$ fixes $\pi'$ pointwise: if $X P_i = P_i X$ for all $i$, then $\pinch_\pi(X) = X$ (\texttt{centerExpectation\_fixes});
\item $\pinch_\pi \circ \pinch_\pi = \pinch_\pi$: the pinching is a projector (\texttt{centerExpectation\_idem});
\item star equivariance: $\pinch_\pi(X)^{*} = \pinch_\pi(X^{*})$ (\texttt{conjTranspose\_centerExpectation}).
\end{enumerate}
\end{proposition}

\begin{proposition}[State preservation; consumes a tracial state]
\label{prop:pinch-state}
If $\rho$ is a state, then $\pinch_\pi(\rho)$ is a state (\texttt{centerExpectation\_isState}).
\end{proposition}

\begin{proposition}[Trace-inner-product geometry; consumes a tracial state]
\label{prop:geometry}
Write $\ip{X}{Y} = \Tr(X^{*}Y)$. For all $X, Y \in \Mn$:
\begin{enumerate}[(i)]
\item selfadjointness: $\ip{\pinch_\pi(X)}{Y} = \ip{X}{\pinch_\pi(Y)}$ (\texttt{trace\_conjTranspose\_centerExpectation\_mul});
\item Pythagoras:
\[
\ip{X}{X} = \ip{\pinch_\pi(X)}{\pinch_\pi(X)} + \ip{X - \pinch_\pi(X)}{X - \pinch_\pi(X)}
\]
(\texttt{trace\_centerExpectation\_pythagoras});
\item squared-$L^2$ contraction, in the partial order of $\C$:
$\ip{\pinch_\pi(X)}{\pinch_\pi(X)} \le \ip{X}{X}$ (\texttt{trace\_centerExpectation\_contraction}).
\end{enumerate}
\end{proposition}

This is the quantum counterpart of the classical $L^2$ energy identity for conditional expectations: the pinching is the orthogonal projection, for the trace inner product, onto the commutant of the context.

\begin{proposition}[Trace compatibility and uniqueness; consumes a tracial state]
\label{prop:unique}
\begin{enumerate}[(i)]
\item For every $X \in \Mn$ and every $C \in \pi'$,
\[
\Tr\bigl(\pinch_\pi(X)\, C\bigr) = \Tr(X C)
\qquad (\texttt{trace\_centerExpectation\_mul\_central}).
\]
\item Conversely, if $Y \in \pi'$ satisfies $\Tr(YC) = \Tr(XC)$ for every $C \in \pi'$, then $Y = \pinch_\pi(X)$ (\texttt{centerExpectation\_unique}).
\end{enumerate}
\end{proposition}

\begin{proof}[Proof sketch]
Item (i) is a cyclic-trace computation using absorption. For (ii), set $D = Y - \pinch_\pi(X)$, observe $D \in \pi'$, test the hypothesis and item (i) against $C = D^{*}$, and conclude $\Tr(D D^{*}) = 0$; faithfulness of the trace forces $D = 0$.
\end{proof}

Item (i) is the defining property of a conditional expectation, and (ii) says the pinching is the only commutant-valued map with that property: existence and uniqueness are both machine-checked.

\begin{proposition}[Compatibility with conditioning; consumes a tracial state]
\label{prop:compat}
Let $P \in \pi'$ (in particular, any central event). For every $\rho \in \Mn$:
\begin{enumerate}[(i)]
\item Born-weight invariance: $\mu_{\pinch_\pi(\rho)}(P) = \mu_\rho(P)$ (\texttt{bornWeight\_centerExpectation});
\item conditioning commutes with the center expectation:
\[
\pinch_\pi\bigl(\lued_P(\rho)\bigr) = \lued_P\bigl(\pinch_\pi(\rho)\bigr)
\qquad (\texttt{centerExpectation\_luedersUpdate}).
\]
\end{enumerate}
\end{proposition}

Item (i) guarantees that the normalizing scalars on both sides of (ii) agree, so the exchange law holds at the level of matrices, degenerate cases included. Conditioning on central events therefore interacts with the center exactly as classical conditional probability interacts with a sub-$\sigma$-algebra: this pair of lemmas is the precise sense in which the center is the classical face of the event algebra.

\section{The expectation functional}
\label{sec:functional}

The fourth module attaches to each state its expectation functional on observables and proves that it is a positive normalized linear functional, strengthening Born-weight positivity from events to arbitrary positive-semidefinite observables.

\begin{definition}[Expectation functional]
\label{def:expectation}
For $\rho, M \in \Mn$, write $\langle M \rangle_\rho \coloneqq \Tr(\rho M)$. Lean name: \texttt{expectation}.
\end{definition}

\begin{proposition}[Positive normalized functional; consumes a tracial state]
\label{prop:functional}
For all $\rho, M, N \in \Mn$ and $c \in \C$:
\begin{enumerate}[(i)]
\item the Born weight is the restriction of the functional to events: $\mu_\rho(P) = \langle P \rangle_\rho$ (\texttt{bornWeight\_eq\_expectation});
\item additivity and homogeneity: $\langle M + N \rangle_\rho = \langle M \rangle_\rho + \langle N \rangle_\rho$ and $\langle cM \rangle_\rho = c \langle M \rangle_\rho$ (\texttt{expectation\_add}, \texttt{expectation\_smul});
\item normalization: $\langle \one \rangle_\rho = 1$ for a state $\rho$ (\texttt{expectation\_one});
\item positivity: if $\rho$ and $M$ are both positive semidefinite, then $0 \le \langle M \rangle_\rho$ in the partial order of $\C$ (\texttt{expectation\_nonneg}).
\end{enumerate}
\end{proposition}

\begin{proof}[Proof sketch]
Positivity no longer has an idempotent to compress with, so the proof decomposes $M$ spectrally: $M = V D V^{*}$ with $V$ unitary and $D$ a nonnegative diagonal, cycles the trace to $\Tr\bigl((V^{*} \rho V) D\bigr)$, and sums termwise products of nonnegative diagonal entries of the sandwiched state against nonnegative eigenvalues. Section~\ref{sec:gaps} explains why this elementary route was chosen over the abstract Gram-factorization route.
\end{proof}

\section{The Tsirelson bound}
\label{sec:tsirelson}

The fifth module is entirely algebra-only: no state and no trace appear anywhere in it. Its content splits into an exact ring identity, two C*-norm lemmas, the abstract bound, and interoperability statements.

\subsection{The ring identity}

\begin{proposition}[CHSH square identity; algebra-only]
\label{prop:ring-identity}
Let $R$ be any ring and $a_0, a_1, b_0, b_1 \in R$ with $a_i^2 = b_j^2 = 1$ and $a_i b_j = b_j a_i$ for all $i, j \in \{0,1\}$. With $S = a_0 b_0 + a_0 b_1 + a_1 b_0 - a_1 b_1$,
\[
S^2 = 4 \cdot 1 - (a_0 a_1 - a_1 a_0)(b_0 b_1 - b_1 b_0)
\qquad (\texttt{chsh\_mul\_self}).
\]
\end{proposition}

The identity involves no star, no order, and no norm: it is the exact algebraic content of the Tsirelson bound, stated in a bare ring. The sign convention is that of commutators $[x,y] = xy - yx$, with the minus sign in front of the product of commutators. As noted in Section~\ref{sec:arch}, this lemma's axiom audit reports only \texttt{[propext, Quot.sound]}. Its proof engineering, a hand-rolled confluent rewrite system standing in for a missing noncommutative tactic, is described in Section~\ref{sec:gaps}.

\subsection{The norm bound}

The analytic half runs in an abstract typeclass context: a normed ring carrying a star-ring structure satisfying the C*-identity $\norm{x^{*}x} = \norm{x}^2$, and nontrivial (so that $\norm{\one} = 1$). In Lean the hypotheses are the typeclasses
\[
[\texttt{NormedRing } A]\ [\texttt{StarRing } A]\ [\texttt{CStarRing } A]\ [\texttt{Nontrivial } A].
\]
Neither completeness nor any order structure on $A$ is assumed; in particular no ambient \texttt{CStarAlgebra} bundle is required.

\begin{proposition}[Norm lemmas; algebra-only]
\label{prop:norm-lemmas}
In the context above:
\begin{enumerate}[(i)]
\item a selfadjoint involution has norm one: if $a^{*} = a$ and $a^2 = \one$, then $\norm{a} = 1$ (\texttt{norm\_eq\_one\_of\_selfAdjoint\_involution});
\item in any normed ring (star and C*-structure are dropped for this item), if $\norm{a_0} = \norm{a_1} = 1$, then $\norm{a_0 a_1 - a_1 a_0} \le 2$ (\texttt{norm\_commutator\_le\_two}).
\end{enumerate}
\end{proposition}

\begin{theorem}[Tsirelson bound, norm form; algebra-only]
\label{thm:tsirelson}
Let $A$ be as above and let $a_0, a_1, b_0, b_1 \in A$ be selfadjoint involutions such that every $a_i$ commutes with every $b_j$. Then
\[
\norm{a_0 b_0 + a_0 b_1 + a_1 b_0 - a_1 b_1} \le 2\sqrt{2}
\qquad (\texttt{tsirelson\_bound}).
\]
\end{theorem}

\begin{proof}[Proof sketch]
The classical square trick, run entirely on the three preceding statements. By Proposition~\ref{prop:ring-identity}, $S^2 = 4\cdot\one - [a_0,a_1][b_0,b_1]$; by Proposition~\ref{prop:norm-lemmas}, each commutator has norm at most $2$, so $\norm{S^2} \le 4 + 4 = 8$. The combination $S$ is selfadjoint, so the C*-identity gives $\norm{S}^2 = \norm{S^{*}S} = \norm{S^2} \le 8$, and $\norm{S} \le \sqrt 8 = 2\sqrt 2$.
\end{proof}

No attainment or tightness statement is formalized anywhere in the module: the development ships the bound as an inequality, and nothing in it asserts that $2\sqrt 2$ is achieved.

\subsection{Relation to Mathlib's CHSH file and the matrix instantiation}

Mathlib already contains an order-theoretic form of the Tsirelson inequality: the file \texttt{Mathlib/Algebra/Star/CHSH} proves, in a star-ordered algebra over the reals, that the CHSH combination satisfies $S \le 2\sqrt{2} \cdot \one$ in the algebra's order, together with the hypothesis bundle \texttt{IsCHSHTuple} packaging selfadjointness, involutivity, and the cross commutations. The present development contributes the complementary \emph{norm-form} statement in the C*-setting, with the underlying algebraic identity isolated as Proposition~\ref{prop:ring-identity}, and connects the two through the shared interface:

\begin{proposition}[Interoperability and instantiation; algebra-only]
\label{prop:tsirelson-instances}
\begin{enumerate}[(i)]
\item The conclusion of Theorem~\ref{thm:tsirelson} holds for any tuple satisfying Mathlib's \texttt{IsCHSHTuple} hypothesis bundle (\texttt{tsirelson\_bound\_of\_isCHSHTuple}).
\item For every $n \neq 0$, the bound holds in $\Mn$ equipped with the $L^2$ operator norm, activated by the scoped instance set \texttt{Matrix.Norms.L2Operator} (\texttt{matrix\_tsirelson\_bound}).
\end{enumerate}
\end{proposition}

Item (ii) returns the abstract statement to the ambient algebra of the event-algebra modules; combined with Proposition~\ref{prop:closure}(ix), any four events with the appropriate commutation pattern yield dichotomic observables to which the bound applies. Because the abstract theorem is typeclass-generic, the matrix case is an instantiation of a few lines rather than a separate proof; the only genuine work is a hand-constructed nontriviality instance for the matrix algebra (Section~\ref{sec:gaps}).

\section{Mathlib gaps and proof engineering}
\label{sec:gaps}

This section reports the friction points and affordances encountered against the pinned Mathlib version, as raw material for library builders and for others formalizing quantum-probabilistic content.

\subsection{Scoped instances}

Three layers of structure needed by this development are scoped. The partial order on $\C$ requires \texttt{open scoped ComplexOrder}; without it, goals mentioning $0 \le z$ or \texttt{Matrix.PosSemidef} over $\C$ fail to elaborate, with instance errors that do not point to the missing scope. The Loewner order on matrices is likewise scoped (\texttt{MatrixOrder}), as is the ordered-ring structure on $\C$ that supplies \texttt{mul\_nonneg} in the complex order. Most consequentially, there is no globally installed norm on the matrix type: the $L^2$ operator norm and its \texttt{NormedRing} and \texttt{CStarRing} instances activate only under \texttt{open scoped Matrix.Norms.L2Operator}, so the matrix instantiation of Theorem~\ref{thm:tsirelson} must live behind that scope. In the same corner, no \texttt{Nontrivial} instance fired for the matrix algebra, and nontriviality (needed for $\norm{\one} = 1$) had to be constructed by hand, entrywise, from $n \neq 0$. The scoping policy is defensible library design, since competing orders and norms exist on the same types; downstream, it is a real usability tax, and the failure modes when a scope is missing are hard to diagnose.

\subsection{Positivity of the trace pairing}

Mathlib provides $\Tr M \ge 0$ for a single positive-semidefinite matrix and covers sandwiches $B A B^{*}$ of positive-semidefinite matrices, yet the bilinear statement, $\Tr(AB) \ge 0$ for positive-semidefinite $A$ and $B$, is absent. For Born weights this is harmless: an event is idempotent, so $\Tr(\rho P) = \Tr(P \rho P)$ by cycling the trace, and the single-matrix lemma applies. For the expectation functional (Proposition~\ref{prop:functional}(iv)) there is no idempotent to compress with, and the natural abstract route failed in an instructive way. Mathlib's Gram-factorization characterization of positivity, \texttt{CStarAlgebra.nonneg\_iff\_eq\_star\_mul\_self} combined with the Loewner-order bridge \texttt{Matrix.PosSemidef.nonneg}, is gated behind the continuous functional calculus, and instance synthesis for two of its classes, \texttt{NonUnitalContinuousFunctionalCalculus} over $\R$ with the \texttt{IsSelfAdjoint} predicate and \texttt{NonnegSpectrumClass} over $\R$, did not resolve for the matrix algebra carrying the Pi topology, even with the order scopes open. (Mathlib's own matrix-order file uses the lemma successfully inside its own section context; the instance chain did not reproduce downstream for us.) The workaround is elementary and robust: decompose spectrally via \texttt{Matrix.IsHermitian.spectral\_theorem}, extract nonnegative eigenvalues with \texttt{PosSemidef.eigenvalues\_nonneg}, cycle the trace, and sum termwise; about a dozen lines.

\subsection{No noncommutative \texttt{linear\_combination}}

The CHSH square identity (Proposition~\ref{prop:ring-identity}) is pure noncommutative-ring algebra with side relations: involutivity of the four letters and the four cross commutations. Lean's \texttt{linear\_combination} tactic is commutative-only, and \texttt{noncomm\_ring} normalizes without consuming hypotheses, so neither closes the goal. The formal proof instead builds a small confluent rewrite system by hand: from each commutation hypothesis it derives the association-compatible, universally quantified form $b\,(a\,x) = a\,(b\,x)$, and from each involution the cancellation form $a\,(a\,x) = x$; running these as a \texttt{simp only} set over right-associated words moves every $b$ past every $a$ and cancels adjacent equal letters, and \texttt{abel} finishes the additive bookkeeping. The system terminates because each rewrite strictly decreases the number of letters out of normal order. A noncommutative analogue of \texttt{linear\_combination}, even one restricted to rewriting modulo hypothesis-supplied commutation and involution relations, would compress this style of proof considerably.

\subsection{Occurrence-sensitive rewriting}

A recurring low-level friction: on goals containing several syntactic occurrences of the same subterm, \texttt{rw} rewrites the first occurrence, which in conditioning goals is typically the normalizing scalar $(\mu_\rho(P))^{-1}$ rather than the intended trace argument. The stable pattern was to isolate each trace identity as a standalone \texttt{have} with explicit arguments to \texttt{trace\_mul\_comm} or \texttt{trace\_mul\_cycle}, and only then rewrite. A minor lexical note in the same spirit: the script capital letter often used for partitions in informal notation is outside Lean's legal identifier alphabet in this toolchain, so partition variables are plain ASCII.

\subsection{Affordances}

Three parts of Mathlib made the development markedly shorter than a bespoke stack would be. First, the \texttt{Matrix.PosSemidef} API is definition-agnostic: sandwich closure in both directions (\texttt{PosSemidef.mul\_mul\_conjTranspose\_same}, \texttt{PosSemidef.conjTranspose\_mul\_mul\_same}), sums (\texttt{posSemidef\_sum}), scalar multiples (\texttt{PosSemidef.smul}, which accepts a complex scalar that is nonnegative in the complex order, so the L\"uders normalization is a single application), diagonal nonnegativity, trace nonnegativity, the zero-trace criterion, and the Gram construction \texttt{posSemidef\_conjTranspose\_mul\_self} covered every need without ever unfolding the definition. Second, trace faithfulness is in the box: \texttt{Matrix.trace\_mul\_conjTranspose\_self\_eq\_zero\_iff} powers the uniqueness of the center expectation, and \texttt{PosSemidef.trace\_eq\_zero\_iff} together with \texttt{PosSemidef.dotProduct\_mulVec\_zero\_iff} powers the support lemma of Proposition~\ref{prop:fixed}(ii). Third, the C*-norm layer, \texttt{CStarRing.norm\_star\_mul\_self}, \texttt{IsSelfAdjoint.norm\_mul\_self}, the norm-one instance for nontrivial C*-rings, and \texttt{norm\_nsmul\_le}, reduces the analytic half of Theorem~\ref{thm:tsirelson} to roughly forty lines over the bare typeclass context, with no completeness assumption anywhere. Mathlib's existing CHSH file was itself an affordance: its \texttt{IsCHSHTuple} bundle fixed the interface against which the norm-form result could be made interoperable from the start.

\section{Related work}
\label{sec:related}

\textbf{The Isabelle/HOL line.} The nearest prior art is the sequence of developments by Echenim and Mhalla. Echenim's Archive of Formal Proofs entry on quantum projective measurements and the CHSH inequality \cite{echenim2021afp} formalizes projective measurements and the CHSH expectation in Isabelle/HOL; the journal treatment by Echenim and Mhalla \cite{echenimmhalla2024} presents a formalization of the CHSH inequality and Tsirelson's upper bound, with the corresponding AFP entry \cite{echenimmhalla2023afp}; a supporting AFP entry formalizes simultaneous diagonalization of pairwise commuting Hermitian matrices \cite{echenim2022afp}. Relative to that line, the present development differs on two axes. Organizationally, it is event-algebra-first: the theory is arranged around the closure algebra of Hermitian idempotents, a distinguished commutative center with an explicit conditional-expectation projector carrying existence, uniqueness, and contractivity, and a lemma-by-lemma, machine-visible split between algebra-only content and content consuming a tracial state; the Isabelle line is organized matrix-first over a concrete density-matrix and $L^2$-norm stack, with the CHSH development as the target. On the prover axis, the ecosystems differ in a way that shows up in proof economy: in Lean/Mathlib the norm-form Tsirelson bound is proved abstractly for any nontrivial unital C*-ring, with the matrix case a short instantiation, and Born-weight positivity and boundedness live directly in the partial order of $\C$ with no bespoke real-part bookkeeping; conversely, the concrete Isabelle stack obtains from its own spectral toolbox some facts (such as Gram factorization of positive matrices) that are currently gated behind unresolved instance chains in Mathlib, as detailed in Section~\ref{sec:gaps}. The CHSH arithmetic itself is common ground; the organization is the delta.

\textbf{Formalized quantum mathematics more broadly.} Bordg, Lachnitt, and He verified quantum algorithms and information-theoretic results in Isabelle/HOL \cite{bordg2021}, with the accompanying Isabelle Marries Dirac library \cite{bordg2020afp}. In Lean/Mathlib, Dupuis, Lewis, and Macbeth developed formalized functional analysis with semilinear maps \cite{dupuis2024}, infrastructure relevant to any future infinite-dimensional extension of the present work. Garrigue and Saikawa formalize typed compositional quantum computation with lenses \cite{garrigue2026}. Reports on individual Mathlib formalization efforts, such as Nash on Engel's theorem \cite{nash2023} and Baanen and collaborators on Dedekind domains and class groups of global fields \cite{baanen2022}, document both the mathematics and the practice of folding formalization-driven contributions back into the library; the library itself is described in the Mathlib community paper \cite{mathlib2020}, and the Lean~4 system in de Moura and Ullrich \cite{lean4}.

\textbf{Foundations context.} The mathematical foundations of quantum probability remain under active discussion in this journal's scope. Reddiger and Poirier examine the one-body Born rule in curved spacetime \cite{reddigerpoirier2025}; Reddiger analyzes the applicability of Kolmogorov's probability theory to quantum phenomena \cite{reddiger2026}; Aharonov and Shushi revisit the derivation of the Born rule \cite{aharonovshushi2024}, and Beck surveys the textbook orthodoxy on measurement and state update \cite{beck2026}. Arsiwalla, Chester, and Kauffman reconstruct expectation values from a pre-quantum operator algebra, treating classical and quantum wave functions as consequences of that operator layer \cite{arsiwalla2024}; Kumar and Pan study network nonlocality \cite{kumarpan2023}, and Held derives the Tsirelson bound from geometric-algebra assumptions \cite{held2021}. Goldstein bootstraps infinite-dimensional state spaces from finite ones, deriving the von Neumann--L\"uders postulate for projective observables \cite{goldstein2026}. The present paper contributes to that discussion at the level of certified bookkeeping: for the finite core, every hypothesis of every standard statement is now explicit and machine-audited, so debates about what follows from what can point to a checked artifact.

\section{Conclusion and outlook}
\label{sec:conclusion}

The finite event-algebra core of quantum probability admits a complete machine-checked treatment at modest cost: 64 lemmas and 8 definitions cover event closure, Born weights with reality and the probabilistic bounds in the strong complex-order form, repeatable L\"uders conditioning with its fixed-point characterization, a conditional expectation onto a distinguished commutative center satisfying the full classical package (projector, state-preserving, selfadjoint, Pythagorean, contractive, trace-compatible, unique, and commuting with conditioning on central events), a positive normalized expectation functional, and the Tsirelson bound in an abstract C*-setting. The axiom audit is part of the source and re-executes on every build; the lemma-level split between algebra-only and state-consuming content gives the development a second, structural payload beyond the individual theorems: it is a machine-visible map of where probability enters the quantum formalism.

Three extensions are natural. First, unsharp measurements: effects ($0 \le E \le \one$ in the Loewner order) and instruments generalize events and the L\"uders map, and the certainty-set and fixed-point interface of Section~\ref{sec:lueders} suggests the shape the generalized statements should take. Second, attainment: complementing Theorem~\ref{thm:tsirelson} with a formalized two-qubit witness achieving $2\sqrt 2$ would close the gap between the inequality and the bound's sharpness, and is well within reach of the matrix instantiation already in place. Third, infinite dimensions: the finite development deliberately avoids completeness, so its Tsirelson theorem already applies to any nontrivial unital C*-ring, yet the event-and-state layer is genuinely finite. Goldstein's bootstrap of infinite-dimensional state spaces from finite ones in this journal \cite{goldstein2026} marks exactly the frontier a formal development would next confront: which parts of the finite bookkeeping survive the passage verbatim, and which require new analysis (normality of states, $\sigma$-additivity, von Neumann algebras in place of matrix algebras). The conditional-expectation package of Section~\ref{sec:center} is the component we expect to generalize most cleanly, since its statements are already phrased against a commutant rather than against matrix entries.

When the bookkeeping of a foundational formalism is mechanical and audited, discussion can concentrate on the questions that are actually open. For finite quantum event algebras, that state has now been reached.

\section*{Declarations}

\textbf{Funding.} No funding was received for this work.

\textbf{Competing interests.} The author declares no competing interests.

\textbf{Author contributions.} Bernhard Mueller is the sole author: he designed and wrote the Lean development, conducted the axiom audit, and wrote the manuscript.

\textbf{Data availability.} No data were generated or analyzed in this work.

\textbf{Code availability.} The complete Lean development \cite{eventalgebra} accompanies the submission as a self-contained, publicly available Lake project archive. It builds free of \texttt{sorry} placeholders with the toolchain \texttt{leanprover/lean4:v4.29.1} and the Mathlib version pinned in the project manifest; the axiom audits reported in the text re-execute as part of the build.

\bibliographystyle{plainnat}
\bibliography{machine_checked_finite_event_algebras}

\end{document}
